\section{Определение и характеристика взлома функции вознаграждения}

\subsection{Ключевые идеи и математическое описание метода}

В статье~\cite{skalse2025defcharact} авторы представляют формальное определение феномена взлома функции вознагражения — ситуации, когда оптимизация прокси-награды приводит к ухудшению результата по истинной функции награды.

Они вводят строгие определения \emph{взламываемости (hackability)}, \emph{невзламываемости (unhackability)} и \emph{упрощения (simplification)} между функциями вознаграждения и показывают условия их существования.


\subsection{Математическое описание метода}

Авторы рассматривают обучение с подкреплением как задачу максимизации ожидаемой награды агента, взаимодействующего со средой, представленной в форме марковского процесса принятия решений.

Формально марковский процесс принятия решений задаётся шестёркой:
\begin{equation}
\begin{aligned}
    &\mathcal{M} = (S, A, T, I, R, \gamma), \\
    &\text{где } S\text{~-- конечное множество состояний среды,} \\
    &A\text{~-- множество действий агента,} \\
    &T : S \times A \to \Delta(S)\text{~-- функция переходов,} \\
    &I \in \Delta(S)\text{~-- распределение начальных состояний,} \\
    &R\text{~-- функция вознаграждения,} \\
    &\gamma \in [0,1]\text{~-- коэффициент дисконтирования.}
\end{aligned}
\end{equation}

Множество $\Delta(X)$ обозначает множество всех вероятностных распределений над множеством $X$.


\subsubsection{Политика и траектории}

\emph{Политика} $\pi : S \to \Delta(A)$ определяет вероятностное распределение действий, выбираемых агентом в каждом состоянии.
Траектория взаимодействия агента со средой описывается последовательностью

\begin{equation}
\begin{aligned}
    \tau &= (s_0, a_0, r_0, s_1, a_1, r_1, \ldots), \\
    &\text{где } s_t\text{~-- состояние в момент времени } t, \\
    &a_t\text{~-- выбранное действие,} \\
    &r_t\text{~-- награда.}
\end{aligned}
\end{equation}

Ожидаемое \emph{возвратное вознаграждение} для политики $\pi$ определяется как:

\begin{equation}
G(\tau) = \sum_{t=0}^{\infty} \gamma^t r_t, \qquad J(\pi) = \mathbb{E}_{\tau \sim \pi}[G(\tau)].
\end{equation}

Значение $J(\pi)$ интерпретируется как ожидаемая дисконтированная награда при следовании политике $\pi$.


\subsubsection{Функция посещаемостей состояний-действий}

Для анализа поведения политики авторы вводят вектор \emph{дисконтированных частот посещений состояний-действий} $F_\pi(s,a)$:

\begin{equation}
F_\pi(s,a) = \mathbb{E}*{\tau \sim \pi}\left[ \sum*{t=0}^{\infty} \gamma^t \mathbf{1}(s_t = s, a_t = a) \right].
\end{equation}

Эта функция описывает, как агент в среднемвыполняет действие $a$ в состоянии $s$,
с учётом дисконтирования будущих шагов.

Тогда значение политики можно выразить как скалярное произведение между функцией награды и функцией посещаемостей:

\begin{equation}
J(\pi) = \langle R, F_\pi \rangle = \sum_{s,a} R(s,a) F_\pi(s,a).
\end{equation}

Это выражение показывает, что функция награды $R$ индуцирует линейное отображение из пространства политик в пространство ожидаемых значений.


\subsubsection{Взламываемость и упрощение награды}

Чтобы формализовать явление взлома функции вознаграждения, авторы вводят понятие \emph{взламываемости} между двумя функциями вознаграждения.


\paragraph{Определение 1. Взламываемость.} Пара функций вознаграждения $R_1, R_2$ называется \emph{взламываемой} относительно множества политик $\Pi$,
если существуют $\pi, \pi' \in \Pi$ такие, что выполняется:

\begin{equation}
J_1(\pi) < J_1(\pi') \quad \text{и} \quad J_2(\pi) > J_2(\pi').
\end{equation}

В противном случае $(R_1, R_2)$ считаются \emph{невзламываемыми}.

Иными словами, улучшение политики по одной награде не должно приводить к ухудшению по другой.

Невзламываемость является симметричным свойством, однако не транзитивным: например, константная функция награды является невзламываемой с любой другой функцией, но не задаёт содержательного сравнения.

\paragraph{Определение 2. Упрощение.} Функция вознаграждения $R_2$ является \emph{упрощением} функции $R_1$ относительно множества политик $\Pi$, если:

\begin{equation}
\begin{aligned}
&J_1(\pi) < J_1(\pi') \Rightarrow J_2(\pi) \le J_2(\pi'), \
J_1(\pi) = J_1(\pi') \Rightarrow J_2(\pi) = J_2(\pi'), \\
&\text{и существует хотя бы одна пара политик} \pi, \pi' \in \Pi, \text{для которых} \\
&J_1(\pi) \ne J_1(\pi'), \text{ но } J_2(\pi) = J_2(\pi').
\end{aligned}
\end{equation}

Если $R_2$ является константой на $\Pi$, то это \emph{тривиальное упрощение}.

Интуитивно, упрощение — это операция, которая \emph{схлопывает различия} между политиками,
заменяя строгие неравенства равенством, но не меняя порядок предпочтений.

\vspace{11pt}
Таким образом, авторы формулируют строгую математическую основу для анализа взаимосвязей между функциями награды,
что позволяет формально описывать и классифицировать случаи, когда прокси-награда безопасно аппроксимирует истинную,
а также ситуации, в которых возникает \emph{reward hacking}.


\subsection{Алгоритм}

Вместо алгоритма в привычном вычислительном смысле, статья формулирует математические условия и процедуру проверки существования невзламываемых и упрощённых функций вознаграждения.
Работа последовательно рассматривает три случая:

\begin{enumerate}
    \item все возможные политики,
    \item конечные множества политик,
    \item частично бесконечные множества без открытых подмножеств.
\end{enumerate}

\subsubsection{Невозможность невзламываемости на множестве всех политик}

\paragraph{Теорема 1.} Если множество политик $\widehat{\Pi}$ содержит открытое подмножество, то любые две нетривиальные функции вознаграждения $R_1$ и $R_2$, невзламываемые на $\widehat{\Pi}$, эквивалентны на $\widehat{\Pi}$:

\begin{equation}
J_1(\pi) = J_1(\pi') \iff J_2(\pi) = J_2(\pi'), \quad \forall \pi, \pi' \in \widehat{\Pi}.
\end{equation}


\paragraph{Следствие 1.} В любом MDP без награды $\mathcal{M} \setminus R$ любые нетривиальные и невзламываемые функции $R_1, R_2$ на множестве всех стационарных политик $\Pi$ эквивалентны.

\vspace{11pt}
Интерпретация: если множество политик имеет непрерывное подмножество в пространстве стратегий, то линейная структура $J(\pi) = \langle R, F_\pi \rangle$ не допускает несовместимых упорядочений политик. Следовательно, все невзламываемые функции награды оказываются эквивалентными по индуцируемому порядку предпочтений.


\subsubsection{Конечные множества политик}

\paragraph{Теорема 2.} Для любого MDP без награды $\mathcal{M} \setminus R$, любого конечного множества политик $\widehat{\Pi}$, содержащего хотя бы две различные политики $\pi, \pi'$, такие что $F(\pi) \ne F(\pi')$, 
и любой функции вознаграждения $R_1$, существует нетривиальная функция $R_2$, такая что $R_1$ и $R_2$ невзламываемы, но не эквивалентны.

\vspace{11pt}
\emph{Схема доказательства}: при непрерывном изменении направления функции вознаграждения $R_1$ в пространстве векторов $F(\pi)$ можно найти промежуточное положение $R_2$, при котором одно из неравенств между значениями политик становится равенством, не изменяя общий порядок предпочтений. Таким образом, $R_2$ является невзламываемой, но нетривиальной функцией вознаграждения.


\subsubsection{Проверка существования упрощения}


\paragraph{Теорема 3.} Пусть $\widehat{\Pi}$ — конечное множество политик, а $R_1$ — функция вознаграждения.  

Тогда процедура проверки существования нетривиального упрощения включает следующие шаги:

\begin{enumerate}
    \item Разделить множество политик $\widehat{\Pi}$ на подмножества $E_1, \dots, E_m$, где
    \begin{equation}
    \pi, \pi' \in E_i \iff J_1(\pi) = J_1(\pi').
    \end{equation}
    \item Для каждого множества $E_i$ выбрать представителя $\pi_i \in E_i$ и построить множество векторов:
    \begin{equation}
    Z_i = \{ F(\pi) - F(\pi_i) : \pi \in E_i \}.
    \end{equation}
    \item Упрощение существует тогда и только тогда, когда
    \begin{equation}
    \dim(Z_1 \cup \dots \cup Z_m) \le \dim(F(\widehat{\Pi})) - 2.
    \end{equation}
\end{enumerate}

Интерпретация: если множество политик $\widehat{\Pi}$ имеет ограниченное разнообразие векторов посещаемостей $F(\pi)$, то оно накладывает недостаточное число независимых ограничений на функцию вознаграждения. В этом случае часть различий между политиками может быть устранена без нарушения порядка предпочтений.


\paragraph{Следствие 3.} Если для всех $\pi, \pi' \in \widehat{\Pi}$ выполняется $J(\pi) \ne J(\pi')$,  
то существует нетривиальное упрощение $R$.

\subsubsection{Бесконечные множества без открытых подмножеств}

Для некоторых бесконечных, но не открытых множеств политик возможно существование нетривиальных упрощений, а для других — нет.  

\vspace{11pt}
Пример с тремя политиками $A, B, C$ и функциями вознаграждения $R_1, R_2$ демонстрирует, что при одной конфигурации $R_2$ упрощает $R_1$, а при другой нарушает порядок предпочтений.

\vspace{11pt}
Таким образом, предложенная процедура задаёт строгие условия существования невзламываемых и упрощённых функций вознаграждения и показывает, что такие случаи ограничены: они возможны лишь при сужении множества политик или контроле процесса оптимизации.



\subsection{Наборы данных}

Работа имеет теоретический характер и не использует внешние наборы данных.  Анализ проводится на основе формального определения марковского процесса принятия решений и рассматривает абстрактные пространства состояний и действий.  

\vspace{11pt}
Для иллюстрации вводится концептуальный пример~-- \emph{Cleaning Robot},  
в котором агент управляет роботом, очищающим три комнаты: чердак, спальню и кухню.  
Этот пример служит демонстрацией того, как частичное или искажённое задание функции вознаграждения может приводить к взлому функции вознаграждения.  

В частности, если истинная функция вознаграждения определяется как
\begin{equation}
r_{\text{true}} = [1, 1, 1],
\end{equation}
а прокси-награда, переданная агенту, имеет вид
\begin{equation}
r_{\text{proxy}} = [1, 0, 0],
\end{equation}
то поведение агента, оптимизирующего $r_{\text{proxy}}$, приводит к снижению значения истинной награды $r_{\text{true}}$.


\subsection{Метрики}

Поскольку работа теоретическая, в ней не используются эмпирические метрики производительности.  
Основным аналитическим инструментом служит \emph{функция возврата} $J(\pi)$~-- ожидаемое дисконтированное вознаграждение при следовании политике $\pi$:
\begin{equation}
J(\pi) = \mathbb{E}_{\tau \sim \pi}\left[\sum_{t=0}^{\infty} \gamma^t r_t \right].
\end{equation}

Сравнение функций вознаграждения $R_1$ и $R_2$ производится через их индуцированные порядки предпочтений над множеством политик:
\begin{equation}
J_1(\pi) < J_1(\pi') \quad \text{и} \quad J_2(\pi) > J_2(\pi').
\end{equation}
Если существует хотя бы одна пара политик, удовлетворяющая этим неравенствам, $(R_1, R_2)$ считаются \emph{взламываемыми}.  
В некоторых случаях дополнительно анализируется \emph{упрощение} (simplification), при котором различие между политиками устраняется без изменения относительного порядка:
\begin{equation}
J_1(\pi) < J_1(\pi') \Rightarrow J_2(\pi) \le J_2(\pi').
\end{equation}


\subsection{Модели}

Работа не опирается на конкретные обучающие модели или нейронные архитектуры.  
Основная модель представлена в виде абстрактного агента, действующего в марковской среде $\mathcal{M} = (S, A, T, I, R, \gamma)$.  

Для различных сценариев рассматриваются:

\begin{itemize}
    \item множество всех политик (включая стохастические), где показано отсутствие нетривиальной невзламываемости;
    \item конечные множества политик, для которых доказано существование нетривиальных невзламываемых пар функций вознаграждения;
    \item частично бесконечные множества без открытых подмножеств, допускающие упрощения в некоторых случаях.
\end{itemize}

Также вводится представление политик через их дисконтированные посещаемости:
\begin{equation}
F_\pi(s,a) = \mathbb{E}_{\tau \sim \pi}\left[\sum_{t=0}^{\infty} \gamma^t \mathbf{1}(s_t = s, a_t = a)\right],
\end{equation}

что позволяет трактовать функцию вознаграждения как линейный функционал на пространстве политик:

\begin{equation}
J(\pi) = \langle R, F_\pi \rangle.
\end{equation}

